% A CVPR-class paper skeleton, proportioned from the corpus.
%
% The comments are word budgets, not suggestions: they are the interquartile
% range of 45 published papers. Delete them as you fill each section in.
%
% Check yourself against the corpus at any point:
%   python3 scripts/audit_paper.py .

\documentclass[10pt,twocolumn,letterpaper]{article}
\usepackage{cvpr}              % or: aaai, neurips, iclr, llncs
\usepackage{graphicx}
\usepackage{booktabs}          % 70% of corpus tables use these rules
\usepackage{amsmath,amssymb}

\begin{document}
\title{A title that names the mechanism, not the outcome}
\maketitle

% ---------------------------------------------------------------- abstract
% Corpus: 149-187 words, 9 sentences. Structure that recurs:
%   1-2  the setting, and the specific failure
%   3    the insight, in one sentence
%   4-6  what we build
%   7-8  the headline numbers, on two benchmarks
%   9    what it means
\begin{abstract}
\end{abstract}

% ------------------------------------------------------------ introduction
% ~900 words, five paragraphs:
%   1  the problem, in the world
%   2  what prior work does and where it breaks -- one paragraph, not three
%   3  the insight, and the figure that shows it (Fig. 1 goes here)
%   4  the method in five sentences
%   5  contributions -- three, each one checkable against a table
\section{Introduction}

% The teaser sits 6.9% into the paper, in 41 of 46 cases as a single graphic,
% single column. Its caption runs ~49 words and states the claim.
\begin{figure}[t]
  \centering
  \includegraphics[width=\linewidth]{figures/teaser.pdf}
  \caption{}
  \label{fig:teaser}
\end{figure}

% ------------------------------------------------------------ related work
% ~700 words. Group by approach, not by year. End each group with the
% sentence that says what it cannot do -- that sentence is what your method
% is measured against, so make it one you can actually beat.
\section{Related Work}

% ------------------------------------------------------------------ method
% ~1400 words, the largest section. Corpus median is 8 equations and 0
% algorithm environments: the maths and one figure carry the method.
% Name this section after your method if it has a name.
\section{Method}
\subsection{Preliminaries}
\subsection{}
\subsection{}

% ------------------------------------------------------------- experiments
% ~1300 words. Order that recurs across the corpus:
%   setup -> implementation details -> main comparison -> ablation
%   -> design study -> cost -> qualitative
% 43 of 46 papers run an ablation. Two benchmarks, not one.
\section{Experiments}
\subsection{Setup}
\subsection{Implementation Details}
\subsection{Comparison with State-of-the-Art}

% Table caption: ~19 words. Say what is being compared and at what setting.
\begin{table}[t]
  \centering
  \caption{}
  \label{tab:main}
  \begin{tabular}{lcc}
    \toprule
    Method & Metric $\uparrow$ & Cost $\downarrow$ \\
    \midrule
    \bottomrule
  \end{tabular}
\end{table}

\subsection{Ablation Study}
\subsection{Analysis}

% ------------------------------------------------------------- conclusion
% ~150 words. What was shown, and the one direction it opens. Not a summary
% of the paper -- the reader just read it.
\section{Conclusion}

{\small\bibliographystyle{ieee_fullname}\bibliography{refs}}
\end{document}
